\subsection{Adapter: Safe Copies and Demand-Driven Recovery}
\label{sec:kv-runtime}

\paragraph{Publish only complete replicas.}
Adapter (4) invokes the Connector to allocate, transfer, and acknowledge
objects selected by the Planner. Peer and fault-independent DRAM paths
commit separately. Publication follows transfer completion and
model/layout/shard validation; a proposed plan is not a readable copy.
Selector state is transferred with the initial layer and must be ready
before selection. A Peer object that is the only surviving copy cannot
be discarded before its DRAM copy commits.

\paragraph{Yield to foreground work.}
On memory pressure, the Adapter withdraws complete Peer footprints in
reverse admission priority, releasing later layers first. Those requests
retain DRAM recovery and compete again in the next window.
Low-priority copy streams are rate-limited (default 20\% of link
capacity), accounting also for the target's HBM bandwidth.

\paragraph{Preserve attention inputs.}
For each retry, preloads overlap earlier-layer computation; newly
revealed misses outrank queued preloads. The runtime's DMA granularity
may limit the model's ideal preemption. A layer executes only when
\emph{all} its required attention inputs are ready; later cold blocks
can arrive in the background or be fetched by subsequent decode.
Tier R rebuilds the history. AegisKV changes placement and loading order,
not selector outputs or known replay tokens. With compatible weights,
layout, and operators it reconstructs the same logical KV state without
an added attention approximation, but does not guarantee identical
future stochastic samples across different floating-point schedules.
